\section{Notation}\label{app:notation}
This appendix collects the notation used throughout the paper, grouped in the order of first
appearance. Unless an explicit comparison across different scales is required, the dependence
of each quantity on $M$ is omitted. The finite-size sign imbalance is denoted by $\rho$ and its
limit by $\rho_\infty$; for the remaining parameters ($r$ and the like), the absence of a
subscript denotes the finite quantity, and the subscript $\infty$ is added when the limit is
meant. The subscripts $c$ and $k$ denote constants that depend only on $k$. The symbols
$\Rightarrow$ and $\xrightarrow p$ denote convergence in distribution and convergence in
probability, respectively. Unless stated otherwise, $k\ge3$ is fixed and $M\to\infty$.

\newcommand{\notationtable}[2]{%
\par\smallskip\noindent\textbf{#1}\par\nopagebreak\smallskip
{\small\noindent\begin{tabular}{@{}p{.22\linewidth}p{.74\linewidth}@{}}
\toprule
Notation & Meaning\\
\midrule
#2
\bottomrule
\end{tabular}}\par\medskip}

\notationtable{Model and peeling}{
$M$ & number of cells, i.e. the number of vertices of the mapping hypergraph\\
$d$ & number of input elements, i.e. the estimation target\\
$k$ & number of cells to which each element is mapped\\
$S_i$ & the $k$-subset chosen by element $i$\\
$F$ & the event that ideal peeling fails, equivalently that the residual 2-core is nonempty\\
$\alpha$ & limit of the load $d/M$, Theorem~\ref{thm:sign}\\
$r$ & window coordinate, $d/M=\alpha_c+r/\sqrt M$; within the critical window $r$ is fixed\\
$n$ & candidate count, used in the test in place of the unknown $d$\\
}

\notationtable{Signs and measurement}{
$s_i$ & sign of element $i$, taking values in $\pm1$\\
$d_+$, $d_-$ & numbers of positive and negative elements\\
$b$ & total sign $d_+-d_-=d\rho$, observable by the protocol\\
$\rho$, $\rho_\infty$ & the sign imbalance $b/d$ and its limit\\
$C_v$ & count of cell $v$, $\sum_{i:v\in S_i}s_i$\\
$T$ & the count energy before peeling, Eq.~\eqref{eq:T}\\
$\widehat d$ & estimate of $d$ obtained from the energy\\
$s$ & standard deviation $\sqrt{2d(d-1)/(M-1)}$; $s(n)$ is the same expression evaluated at a candidate $n$\\
$X$ & standardized measurement $(\widehat d-d)/s$\\
$X_{\rm one}$ & the standardized measurement when all signs are equal\\
$\gamma$ & sign factor $(d\rho^2-1)/(d-1)$, with $\gamma\to\rho_\infty^2$\\
$V$ & sign residual $X-\gamma X_{\rm one}$\\
$c_F$ & increment in the expected overlap of two elements under failure selection relative to the unconditional case, Eq.~\eqref{eq:projection}\\
$q_u$ & the $u$ lower quantile of $\widehat d/d$ under failure\\
}

\notationtable{Critical constants}{
$x_c$ & the unique positive root of the equation $e^{x}-1=(k-1)x$\\
$q_c$, $a_c$ & $1-e^{-x_c}$ and $q_c^{k-1}$\\
$\mu_c$, $\alpha_c$ & critical mean degree $x_c/a_c$ and critical load $\mu_c/k$\\
$\beta_k$ & slope of the score with respect to the window coordinate, $-ka_ce^{-x_c}$\\
$\sigma_k^2$ & limit variance of the score, equal to the initial part $\sigma_0^2$ plus the dynamic part $\tau_k^2$\\
$\tau_k^2$ & variance of the dynamic part of the score, Eq.~\eqref{eq:tau}\\
$\eta_k$ & limit correlation coefficient of the measurement and the score when all signs are equal, Eq.~\eqref{eq:eta}\\
$w_k$ & window width $\sigma_k/|\beta_k|$\\
$C_k$ & coefficient of the finite-size threshold shift, Section~\ref{sec:rules}\\
}

\notationtable{Joint limit and conditional law}{
$R$ & precritical score, depending only on the unsigned mapping, Eq.~\eqref{eq:score}\\
$z_r$ & standardized window coordinate $-\beta_kr/\sigma_k$\\
$c$ & correlation coefficient of the joint limit, $\eta_k\rho_\infty^2$\\
$\Phi$, $\phi$ & standard normal distribution function and density\\
$\Phi_2(\cdot,\cdot;c)$ & standard bivariate normal distribution function with correlation $c$\\
$A(z)$ & $\phi(z)/\Phi(z)$\\
$f_{\rm sel}$, $q_{\rm sel}$ & density and quantile of the selection law\\
$\mathcal G$ & unsigned Plain hypergraph with element labels\\
$\mathcal S$ & event that depends only on the unsigned mapping and is invariant under permutation of the element labels\\
}

\notationtable{Upper confidence bound and implementation rule}{
$u$ & nominal lower-tail level\\
$\mathcal N(b)$ & candidate set $\{n\ge\max(|b|,2):n\equiv b\pmod 2\}$\\
$z(n)$ & standardized window coordinate corresponding to a candidate $n$\\
$t(n,b)$ & acceptance threshold for a candidate $n$\\
$\mathcal C$ & acceptance set; $U=\sup\mathcal C$ is the upper confidence bound\\
$\kappa(u)$ & chi-square shape correction\\
}

\notationtable{Load axis}{
$P_2$ & event that the actual core is exactly two repeated edges\\
$\lambda$, $\lambda_n$ & expected number of repeated edges $\binom d2/\binom Mk$, and its value at a candidate $n$\\
$\delta_k$, $b_k$ & exponents of the error terms in Proposition~\ref{prop:lowload} and Proposition~\ref{prop:mixture}\\
$\mathcal A$ & macroscopic failure event in which the number of edges of the actual 2-core exceeds $\varepsilon_0M$\\
$H_n$, $T_n$, $w_n$ & mixture distribution, acceptance threshold, and macroscopic-layer weight of Proposition~\ref{prop:axis}\\
}

\notationtable{Experimental coordinates}{
$z_0$ & theorem coordinate $(d/M-\alpha_c)\sqrt M/w_k$\\
$z_{\rm sh}$ & shifted coordinate $z_0+C_kM^{-1/6}$\\
$z_{\rm emp}$ & empirical coordinate recovered from the measured failure rate\\
$\widehat p_F$ & measured failure rate\\
$\sigma$ & product of the signs of the two edges of a repeated edge\\
}

\notationtable{Notation in the proofs}{
$D_v$ & initial degree of cell $v$\\
$\mu$ & finite-size mean degree $kd/M$\\
$G(f)$ & centered empirical functional of the initial degree\\
$e$, $h_c$ & influence functions of the initial parts of the measurement and the score\\
$\Gamma(f,g)$ & limit covariance at a fixed total number of endpoints, Eq.~\eqref{eq:bridge}\\
$W$ & dynamic innovation of the score\\
$\Lambda_k$ & limit mean of the number of bad edges in the configuration model\\
$H_{ij}$, $m_{ij}$ & overlap $|S_i\cap S_j|$ of the mappings of two elements\\
$\Omega$, $\Omega_2$, $\ell_i$, $\Omega_\ell$ & total overlap, sum of squares, weighted degree, and sum of squared weighted degrees of the overlap graph\\
$\theta$ & number of steps until the number of degree-one vertices first becomes zero\\
$\vartheta$ & error scale $\sqrt M(\log M)^6$ in the uniform concentration before stopping\\
}
